\documentclass[12pt]{article}

% Layout.
\usepackage[top=1in, bottom=0.75in, left=1in, right=1in, headheight=1in, headsep=6pt]{geometry}

% Fonts.
\usepackage{mathptmx}
\usepackage[scaled=0.86]{helvet}
\renewcommand{\emph}[1]{\textsf{\textbf{#1}}}

% TiKZ.
\usepackage{tikz, pgfplots}
\usetikzlibrary{calc}
\pgfplotsset{compat = newest}
 
\pgfplotsset{my style/.append style={axis x line=middle, axis y line=
middle, xlabel={$x$}, ylabel={$y$}, axis equal }}

% Misc packages.
\usepackage{amsmath,amssymb,latexsym}
\usepackage{graphicx}
\usepackage{array}
\usepackage{xcolor}
\usepackage{multicol}

% Commands to set various header/footer components.
\makeatletter
\def\doctitle#1{\gdef\@doctitle{#1}}
\doctitle{Use {\tt\textbackslash doctitle\{MY LABEL\}}.}
\def\docdate#1{\gdef\@docdate{#1}}
\docdate{Use {\tt\textbackslash docdate\{MY DATE\}}.}
\def\doccourse#1{\gdef\@doccourse{#1}}
\let\@doccourse\@empty
\def\docscoring#1{\gdef\@docscoring{#1}}
\let\@docscoring\@empty
\def\docversion#1{\gdef\@docversion{#1}}
\let\@docversion\@empty
\makeatother

% Headers and footers layout.
\makeatletter
\usepackage{fancyhdr}
\pagestyle{fancy}
\fancyhf{} % Clears all headers/footers.
\lhead{\baselineskip 30pt
%\emph{\@doctitle\hfill\@docdate}
\emph{\@docdate\hfill\@doctitle}
\ifnum \value{page} > 1\relax\else\\
\emph{Name: \rule{3.5in}{1pt}\hfill \@docscoring}\fi}
\rfoot{\emph{\@docversion}}
\lfoot{\emph{\@doccourse}}
\cfoot{\emph{\thepage}}
\renewcommand{\headrulewidth}{0pt}%
\makeatother

% Paragraph spacing
\parindent 0pt
\parskip 6pt plus 1pt

% A problem is a section-like command. Use \problem{5} to
% start a problem worth 5 points.
\newcounter{probcount}
\newcounter{subprobcount}
\setcounter{probcount}{0}
\newcommand{\problem}[1]{%
\par
\addvspace{4pt}%
\setcounter{subprobcount}{0}%
\stepcounter{probcount}%
\makebox[0pt][r]{\emph{\arabic{probcount}.}\hskip1ex}\emph{[#1 points]}\hskip1ex}
\newcommand{\thesubproblem}{\emph{\alph{subprobcount}.}}

% Subproblems are an enumerate-like environment with a consistent
% numbering scheme. 
% Use \begin{subproblems}\item...\item...\end{subproblems}
\newenvironment{subproblems}{%
\begin{enumerate}%
\setcounter{enumi}{\value{subprobcount}}%
\renewcommand{\theenumi}{\emph{\alph{enumi}}}}%
{\setcounter{subprobcount}{\value{enumi}}\end{enumerate}}

% Blanks for answers in normal and math mode.
\newcommand{\blank}[1]{\rule{#1}{0.75pt}}
\newcommand{\mblank}[1]{\underline{\hspace{#1}}}
\def\emptybox(#1,#2){\framebox{\parbox[c][#2]{#1}{\rule{0pt}{0pt}}}}

% Misc.
\renewcommand{\d}{\displaystyle}
\newcommand{\ds}{\displaystyle}
\def\bc{\begin{center}}
\def\ec{\end{center}}
\def\be{\begin{enumerate}}
\def\ee{\end{enumerate}}


\doctitle{Math 251: Quiz 3}
\docdate{11 September, 2025}
\doccourse{UAF Calculus I}
\docversion{v-1}
\docscoring{\blank{0.8in} / 25}
\begin{document}
\textbf{Please circle your instructor's name:} \hfill Kevin Meek \hfill James Gossell \hfill Margaret Short \\

There are 25 points possible on this quiz. No aids (book, calculator, etc.)
are permitted.  {\bf Show all work for full credit.}

\problem{12} Consider the graph of the function $f(x)$ below.

\begin{figure}[h]
\begin{center}
{\includegraphics[width=\linewidth,height=2.25in,keepaspectratio]{plot1.pdf}}
\end{center}
\end{figure}

Use the graph of $f(x)$ to answer each question below. If the limit is infinite, indicate
that with $\infty$ or $-\infty$. If the value does not exist or is undefined, write DNE.

\vspace{.2in}
\begin{tabular}{lll}
(a)  $\displaystyle{\lim_{x \rightarrow 0^-} f(x)} = $  \hspace{21mm} &  
(b) $\displaystyle{\lim_{x \rightarrow 0^+} f(x)} = $ \hspace{21mm} &  
(c) $\displaystyle{\lim_{x \rightarrow 0} f(x)} =  $
\\ \\ \\ \\
(d)  $\displaystyle{\lim_{x \rightarrow 2^-} f(x)} = $  \hspace{21mm} &  
(e) $\displaystyle{\lim_{x \rightarrow 2^+} f(x)} = $ \hspace{21mm} &  
(f) $\displaystyle{\lim_{x \rightarrow 2} f(x)} =  $
\\ \\ \\ \\
(g) $f(0) = $ & 
(h) $f(1) = $ & 
(i) $f(2) = $
\end{tabular}

\vspace{.3in}
\,\, (h) Indicate {\bf{all}} $x$-values for which the function $f(x)$ is {\bf{not continuous}}:

\vspace{.4in}\hspace{.5in} \rule{3.8in}{.1mm}

\pagebreak

\problem{9} Evaluate the following limits. (It is possible that a limit is infinite or
doesn't exist.) Justify your answers. Be sure to use
correct notation for limits in order to receive full credit.
\begin{enumerate}
\item $\displaystyle{\lim_{x \rightarrow -1} \,\, x^2 -3x + \frac{1}{x}}$
\vspace{2.1in}
\item $\displaystyle{\lim_{h \rightarrow 2} \,\, \frac{h^2 -h - 2}{3h^2-9h+6}}$
\vspace{2.1in}
\item $\displaystyle{\lim_{\theta \rightarrow 0} \,\,\frac{ \sin^2 (\theta)}{\theta \cos(\theta)}}$ 
\hspace{1in}
(Hint: $\displaystyle{\lim_{\theta \rightarrow 0} \frac{\sin(\theta)}{\theta} = 1}$.)
\end{enumerate}
\vspace{2.1in}

\pagebreak

\problem{4} Determine whether or not the given function is continuous at
$x = 1$. {\bf{Justify your answer using proper limit notation.}}

$$f(x) = \left\{ \begin{array}{ll}
x^2-2x+3 & {\mbox{if }} x < 1 \\
\frac{5x-29}{2x-10} & {\mbox{if }} x \ge 1 \end{array} \right. $$


\vspace{55mm}
\problem{2} BONUS: Does the equation $x+2x^2=8\sqrt{x}$ have a solution for some
$x$ in the interval [1,9]? Justify your answer. (Hint: Consider the function $f(x) = x+2x^2-8\sqrt{x}$.)


	
\end{document}
%%% Local Variables:
%%% mode: LaTeX
%%% TeX-master: t
%%% End:
